\begin{table}[H]
\centering
\caption{Kebutuhan Fungsional (FR) Sistem}\label{tbl:kebutuhan-fungsional}
% Menambah sedikit jarak antar baris agar tidak terlalu padat
\renewcommand{\arraystretch}{1.3}
\begin{tabular}{|l|p{0.85\textwidth}|}
\hline
\textbf{Kode} & \textbf{Deskripsi Kebutuhan Fungsional} \\
\hline

% Bagian 1
\multicolumn{2}{|l|}{\textbf{Interaksi Pengguna (Siswa)}} \\
\hline
FR-01 & Sistem harus menyediakan mekanisme otentikasi bagi siswa untuk mengakses lingkungan pembelajaran personal mereka. \\
\hline
FR-02 & Sistem harus menyajikan materi pembelajaran yang telah disesuaikan tingkat kesulitannya dengan kemahiran siswa saat ini. \\
\hline
FR-03 & Sistem harus memfasilitasi \textbf{interaksi pembelajaran antar-pengguna} yang memungkinkan terjadinya pertukaran informasi atau kolaborasi akademik di dalam platform. \\
\hline
FR-04 & Sistem harus menyediakan informasi mengenai konteks sosial pembelajaran (misalnya, aktivitas komunitas atau capaian kolektif) untuk membangun kesadaran sosial (\textit{social awareness}). \\
\hline

% Bagian 2
\multicolumn{2}{|l|}{\textbf{Logika Adaptasi \& Pemrosesan Data (Sistem)}} \\
\hline
FR-05 & Sistem harus memiliki algoritma untuk mengestimasi tingkat kemahiran pengetahuan individu siswa secara dinamis berdasarkan riwayat aktivitasnya. \\
\hline
FR-06 & Sistem harus mampu \textbf{mengkuantifikasi data interaksi sosial} (dari FR-03) menjadi parameter numerik yang dapat diproses lebih lanjut (misalnya, skor keterlibatan atau skor kinerja kolaboratif). \\
\hline
FR-07 & Sistem harus memiliki mekanisme untuk menggabungkan parameter kemahiran individu (FR-05) dan parameter sosial (FR-06) dalam satu model komputasi terpadu. \\
\hline
FR-08 & Sistem harus menggunakan hasil penggabungan model (FR-07) untuk menentukan materi atau aktivitas berikutnya yang paling optimal bagi siswa. \\
\hline
FR-09 & Sistem harus mampu merekomendasikan \textbf{intervensi sosial spesifik} (misalnya, merekomendasikan orang untuk berinteraksi atau topik untuk didiskusikan) guna mencegah isolasi pembelajaran. \\
\hline
\end{tabular}
\end{table}